\section{Déformation positive et dualité réticulaire du registre engendré}
\label{rec:positiva}

\subsection{Direction de couplage et flux unimodulaire}

L’extraction du registre, sa lecture rationnelle et les projecteurs d’incidence ont été construits dans les sections précédentes. On conserve la base mille, l’origine de phase, la frontière de retenue et le système de coordonnées des douze positions. Dans ce système de coordonnées,
\[
 K=(234,543,140,729,659,824,621,058,914,794,146,601),
 \qquad u=P_3P_{11}K,
\]
\[
 P_{11}=I-\frac1{12}\mathbf1\mathbf1^{\mathsf T},\qquad
 u\perp\mathbf1,\qquad
 \|u\|^2=\frac{6638585+2275584\sqrt5}{20}>0.
\]
La définition finie de \(P_3\), l’ordre de ses classes et l’évaluation exacte de cette norme figurent dans le lemme~\ref{lem:k-sector-no-nulo}. La réalisation qui suit utilise cet objet engendré. Nous écrivons
\begin{equation}
 \rho=\frac{\pi_{\HMT}}{\varphi_{\HMT}^2}>0,\qquad
 h_i=(\log\rho)\frac{u_i}{\|u\|},\qquad
 r_i(t)=e^{th_i},\qquad g_t=\operatorname{diag}(r_i(t)).
 \label{rec:eq:flujo}
\end{equation}
Les coordonnées régionales sont évaluées avant cette composition. Le paramètre \(t\in\RR\) indexe une déformation mathématique ; sa définition ne l’identifie ni à l’horloge TPK ni à une profondeur de raffinement.


\begin{proposition}[Compensation d’échelle]
\label{rec:prop:flujo}
Le flux vérifie
\[
 g_tg_s=g_{t+s},\qquad g_0=I,\qquad g_t^{-1}=g_{-t},\qquad
 \det g_t=\prod_i r_i(t)=1.
\]
\end{proposition}
\begin{proof}
L’orthogonalité au mode uniforme donne \(\sum_i h_i=0\). La multiplication des exponentielles prouve la loi de groupe et l’inversion. Le déterminant est \(\exp(t\sum_i h_i)=1\). Chaque facteur diagonal est positif.
\end{proof}

La conservation est multiplicative et globale : chaque rayon peut varier. En coordonnées logarithmiques, la somme nulle est conservée. C'est cette distinction qui permettra d'interpréter géométriquement le projecteur transversal de rang dix.

\subsection{Réciprocité exacte du double cercle
}
\label{rec:doble-circulo}

Considérons le système de coordonnées complexe et son inverse
\[
 z(s)=\frac{s-1}{s},\qquad s(z)=\frac1{1-z}.
\]
Pour \(r>0\), \(z=re^{i\theta}\ne1\), nous définissons \(s(r,\theta)=(1-re^{i\theta})^{-1}\).

\begin{proposition}[Conjugaison du miroir spectral
]
\label{rec:prop:espejo}
On a
\begin{align}
 s(r^{-1},\theta)&=1-\overline{s(r,\theta)},\label{rec:eq:espejo}\\
 \Re s(r,\theta)-\frac12&=
 \frac{1-r^2}{2|1-re^{i\theta}|^2}.\label{rec:eq:lado}
\end{align}
Le cercle unité, sauf le point exclu \(z=1\), correspond à la droite \(\Re s=1/2\).
\end{proposition}
\begin{proof}
L’inversion radiale à angle fixe est \(z'=1/\overline z\). Par conséquent,

\[
 \frac1{1-z'}=\frac{\overline z}{\overline z-1}
 =1-\frac1{1-\overline z}.
\]
De plus, \((1-z)^{-1}=(1-\overline z)/|1-z|^2\). Soustraire \(1/2\) de sa partie réelle produit \eqref{rec:eq:lado}.
\end{proof}

La relation avec les rayons de \eqref{rec:eq:flujo} est explicite : \(r_i(-t)=r_i(t)^{-1}\). L’égalité des produits conserve le volume, tandis que la proposition identifie une involution entre systèmes de coordonnées complexes. La localisation des zéros d’une fonction requiert les propriétés de cette fonction ; elle n’est pas une conséquence du produit unitaire des rayons.

Il convient de distinguer trois réalisations circulaires du corpus. Le système de coordonnées spectral envoie \(s=1/2+i\gamma\) sur \(z_\gamma=(\gamma+i/2)/(\gamma-i/2)\). Le cercle dual de la mémoire entière envoie un entier \(n\) sur le caractère \(z^n\). La compression nonadique envoie \(z\) sur \((8+z)/9\), cercle de centre \(8/9\) et de rayon \(1/9\), tangent au cercle unité en \(1\). Le rapport modal \(5/9\) appartient, en revanche, aux valeurs propres \(9,5\) du bloc APP \(\left(\begin{smallmatrix}7&2\\2&7\end{smallmatrix}\right)\). Leurs fonctions restent différenciées lorsqu'on les compose.

\subsection{Strates positroïdales et invariants équilibrés}

Soit \(X\) une matrice réelle de rang \(k\), à douze colonnes, et soit \(\Delta_I(X)\) son mineur maximal de colonnes \(I\). L'action à droite satisfait
\begin{equation}
 \Delta_I(Xg_t)=e^{t\sum_{i\in I}h_i}\Delta_I(X).
 \label{rec:eq:plucker}
\end{equation}
La positivité diagonale conserve tous les signes et tous les mineurs nuls. Ainsi, pour \(1\le k\le11\), une trajectoire commencée dans \(\operatorname{Gr}_{\ge0}(k,12)\) demeure dans la même strate positroïdale. On utilise ici la stratification par mineurs comme réalisation ultérieure ; voir la construction générale de Postnikov~\cite{rec:postnikov}.

\begin{proposition}[Dualité positive de rangs complémentaires
]
Soient \(A_{\pm}=\operatorname{diag}(1,-1,1,-1,\ldots)\) et
\(\mathfrak D(V)=V^\perp A_{\pm}\), avec l’orientation positive de leurs mineurs complémentaires. Alors

\[
 \mathfrak D(Vg_t)=\mathfrak D(V)g_{-t}.
\]
\end{proposition}
\begin{proof}
Comme \(g_t\) est diagonal et auto-adjoint,
\((Vg_t)^\perp=V^\perp g_t^{-1}\). La matrice \(A_{\pm}\) commute avec \(g_t\). Pour le signe des mineurs, le complément orthogonal introduit, à un signe global près, \((-1)^{\sum_{j\in J}j}\). La multiplication des colonnes par \((-1)^{j-1}\) laisse un signe indépendant de \(J\) ; on l’élimine par une unique orientation globale. Il reste des mineurs complémentaires non négatifs.
\end{proof}

Si \(\mathbf1_I+\mathbf1_J=\mathbf1_L+\mathbf1_M\) et si le dénominateur est non nul, l’annulation coordonnée par coordonnée des exposants prouve

\begin{equation}
 \frac{\Delta_I(Xg_t)\Delta_J(Xg_t)}{\Delta_L(Xg_t)\Delta_M(Xg_t)}
 =\frac{\Delta_I(X)\Delta_J(X)}{\Delta_L(X)\Delta_M(X)}.
 \label{rec:eq:cocientes}
\end{equation}
En rang deux, on récupère les quotients de quatre indices et leurs relations d'échange. En revanche, le produit de mineurs complémentaires des deux trajectoires appariées reçoit deux fois le même facteur et peut varier. Le repère circulant \(O_K\), inversible pour la récupération, n'acquiert pas non plus par cette action la propriété d'être totalement non négatif : ses mineurs conservent les signes établis dans sa propre réalisation.

\subsection{Covolume, réseau dual et symétrie ternaire
}

Soit \(\Lambda\) le réseau de Leech marqué construit précédemment, autodual et de covolume un dans la métrique de ses douze plans orthogonaux de type \(A_2\). L’isométrie marquée

\[
 \sigma=\bigoplus_{i=1}^{12}P_i\mathsf C P_i^{-1},\qquad
 \mathsf C^2+\mathsf C+I=0,
\]
préserve le recollement et \(\Lambda\), est d’ordre trois et ne possède aucun vecteur fixe non nul. Les \(P_i\) sont des repères des plans, distincts des projecteurs \(P_3,P_{11}\). Nous définissons

\begin{equation}
 G_t=\bigoplus_{i=1}^{12}r_i(t)I_{A_{2,i}},\qquad
 \Lambda_t=G_t\Lambda.
 \label{rec:eq:red}
\end{equation}

\begin{theorem}[Déformation réticulaire réciproque]
\label{rec:thm:red}
Pour tout \(t\in\RR\),
\[
 \operatorname{covol}(\Lambda_t)=1,\qquad
 \Lambda_t^*=\Lambda_{-t},\qquad \sigma\Lambda_t=\Lambda_t.
\]
Toute isométrie de \(\Lambda\) qui commute avec \(G_t\) préserve également \(\Lambda_t\).
\end{theorem}
\begin{proof}
Chaque rayon agit sur deux coordonnées ; ainsi,
\(\det G_t=\prod_i r_i(t)^2=1\). Pour un opérateur inversible \(G\), un vecteur \(y\) appartient à \((G\Lambda)^*\) exactement lorsque \(G^*y\in\Lambda^*\). C’est pourquoi
\((G\Lambda)^*=G^{-*}\Lambda^*\). En appliquant l’autodualité, l’auto-adjonction et l’inversion, on obtient \(\Lambda_t^*=G_{-t}\Lambda\). Enfin, \(G_t\) est scalaire dans chaque plan et commute avec \(\sigma\), d’où
\(\sigma G_t\Lambda=G_t\sigma\Lambda=G_t\Lambda\). Le même calcul vaut pour les autres isométries commutantes. L’opérateur \(\sigma\) ne change pas, et conserve donc son ordre et son absence de vecteurs fixes.
\end{proof}

Le point \(t=0\) retrouve le réseau initial avec son intégralité, sa parité, son autodualité et sa réalisation exceptionnelle. La collection complète est une collection de réseaux euclidiens unimodulaires au sens de covolume un ; en dehors du point initial, on ne présuppose ni intégralité ni réalisation de Moonshine. Le théorème identifie exactement la symétrie ternaire qui persiste effectivement et la dualité qui relie les deux sens de déformation.

\begin{figure}[H]
\centering
\begin{tikzpicture}[node distance=3mm,>=Latex,every node/.style={font=\fontsize{8.5}{10}\selectfont,align=center},box/.style={draw=ink,rounded corners=2pt,text width=0.16\linewidth,inner sep=3pt,minimum height=25mm}]
\node[box] (a) {APP--TRIT--TPK : état enrichi et continu conjoint};
\node[box,right=of a] (b) {Registre \(K\), représentations régionales et archive bilatérale};
\node[box,right=of b] (c) {Direction récupérable \(u_K\) et flux positif de déterminant un};
\node[box,right=of c] (d) {Rayons conjugués, réseaux duaux et systèmes de coordonnées hyperboliques};
\node[box,right=of d] (e) {Réciprocité spectrale et récupération orientée de l'action};
\draw[->] (a)--(b);\draw[->](b)--(c);\draw[->](c)--(d);\draw[->](d)--(e);
\end{tikzpicture}
\caption{Dépendances de l’extension. Les flèches représentent des compositions spécifiées dans les démonstrations ; la récupération conserve les systèmes de coordonnées et l’orientation.}
\label{fig:rec-causal}
\end{figure}
